% !TeX program = pdflatex
% AAWM Paper A — arXiv 格式（I-3 同款骨架）
\documentclass[11pt]{article}
\usepackage[utf8]{inputenc}
\usepackage[T1]{fontenc}
\usepackage{hyperref}
\usepackage{url}
\usepackage{booktabs}
\usepackage{multirow}
\usepackage{amsfonts}
\usepackage{amsmath}
\usepackage{graphicx}
\usepackage{microtype}
\newcommand{\green}{\mathrm{green}}
\title{AAWM: Agent-Level Watermarking for Black-Box LLM Outputs}
\author{Yin Cui}
\date{September 2026}

\begin{document}
\maketitle

\begin{abstract}
Model-level watermarks require access to LLM internals; cryptographic agent
identity proves who an agent is but not which text it produced. We present
AAWM, an agent-level watermark that operates entirely in the agent
orchestration layer: any black-box LLM output can be watermarked post-hoc,
and leaks can be traced to a specific user among 65,536. AAWM combines two
channels---a Merkle-HMAC paragraph-binding channel for tamper evidence and
a keyed greenlist statistical channel carrying a 16-bit user identifier via
synonym choice---with a capacity-adaptive codec, a registry soft-matcher,
and an abstain protocol that prefers ``unknown'' to a wrong attribution.
Key-property measurements on 1,020 real LLM-generated deliverables
(consulting reports, legal memos, research sections; Chinese and English)
show: on unedited leaks both existence detection and user
attribution succeed on 100\% of 100 real LLM deliverables (Chinese,
1,270 characters median), embedding
raises perplexity by 9.3\% on average (proxy metric), and---honestly---real
LLM paraphrase at 15\% strength collapses both detection and attribution,
defining a clear capability boundary. All artifacts, including derived-salt
recovery after metadata loss, are open source.
\end{abstract}

\section{Introduction}
When a confidential LLM-produced deliverable leaks, only two
attribution technologies are available, and each leaves a gap. \emph{Model-level watermarks}
(KGW-style greenlists \cite{kgw}, tournament sampling \cite{synthid}) bind the mark to the LLM's
sampling process: they require logits or a modified generator, so they
cannot mark output from closed APIs, and they prove ``this text came from
model $M$''---not \emph{which user} received it. \emph{Cryptographic agent
identity} \cite{agentpin} (signed agent cards, DID-based IDs) proves who an agent is, but a
signature on the agent is not a mark on the text: anyone can copy the text
and strip the signature.

We present \textbf{AAWM}, an \emph{agent-level} watermark that closes this
gap. The agent rewrites its own output post-hoc, choosing synonyms according
to a key derived from the agent instance's master key; a verifier with the
same key reads the mark out of the text like an acrostic---including
\emph{which of up to 65,536 users} received this copy. The LLM is a pure
black box; nothing about the model or the API changes.

\paragraph{Contributions.}
\begin{itemize}
\item \textbf{Agent-level positioning} with a concrete threat model
(\S\ref{sec:design}): the watermark lives in the orchestration layer, works
with any black-box LLM, and binds output to agent-instance and user.
\item \textbf{Dual-channel design}: a Merkle-HMAC paragraph-binding channel
(tamper evidence, localization) plus a keyed-greenlist statistical channel
carrying a capacity-adaptive 16-bit user identifier with redundancy
layouts and registry soft-matching (\S\ref{sec:design}).
\item \textbf{An abstain protocol} that outputs ``unknown'' rather than a
wrong user when attribution confidence is insufficient, with an evidence
mask that distinguishes ``no evidence'' from ``contradicting evidence''
(\S\ref{sec:design}).
\item \textbf{Operational resilience}: session salts are derived from the
master key and the embedding-invariant canonicalization of the text, so
tracing survives total metadata loss (\S\ref{sec:design}).
\item \textbf{Evaluation on real LLM deliverables}---1,020 documents from a
32B LLM across three genres in two languages---with honest boundaries:
100\% detection and exact attribution on unedited leaks
($n{=}100$), $+9.3\%$ perplexity (proxy) embedding cost, and complete
collapse under 15\%-strength LLM paraphrase, which we report as a
capability boundary rather than hide (\S\ref{sec:robustness},
\S\ref{sec:quality}).
\end{itemize}

\section{System Design}
\label{sec:design}

\subsection{Threat model and positioning}
The embedding party is an agent platform that serves many users from one
or more black-box LLMs. The verifying party is an auditor holding the same
master key. Attackers may copy, excerpt, reorder, or lightly edit leaked
text; they may also rephrase it with an LLM. We do not claim resistance to
determined semantic rephrasing (\S\ref{sec:robustness}) or key compromise;
the target is \emph{internal leak tracing}, where leaked documents are
typically reposted unedited or with structural edits.

\subsection{Channel B: keyed greenlist statistics with user coding}
Let $K$ be the master key and $s$ a session salt. A word predicate
$\green_K(w)$ (top bit of $\mathrm{HMAC}(K_{\green}, w)$) and a band map
$b_K(w)$ (mod 16) are derived per salt. A curated synonym dictionary
(136 zero-cost Chinese groups; 133 English groups) partitions into bands;
groups whose members are monochromatic under the keyed predicate are
dropped per salt. A 16-bit user id is encoded as band target colors; the
embedder rewrites every dictionary-word occurrence toward its band's target
color, preserving casing, rejecting boundary-unsafe and adjacent-duplicate
choices. The detector recomputes per-band counts with a per-word vote cap
(c=3) to defeat correlated-vote inflation, thresholds $\sum_b |z_b|$ for
existence, and decodes the uid from band signs; a registry of candidates is
scored by soft dot-product matching, and low-margin or colliding candidates
abstain. Capacity adapts to the text (active bands $k$), with optional
$r$-fold redundant layouts (same bit across $r$ interleaved bands) for
crop resistance; uids are masked to the $k$-bit space and matched against
the registry accordingly.

\subsection{Channel A: paragraph binding}
A Merkle tree over HMAC'd normalized paragraphs, with the uid as associated
data, gives tamper evidence: any edit, insertion, deletion, or reordering
fails root verification and localizes to paragraphs. Channel A is brittle
by design (a paraphrase breaks it); Channel B is robust by design; the two
together separate ``was it edited'' (A) from ``who leaked it'' (B).

\subsection{Abstain protocol and evidence mask}
Attribution outputs (uid, user, confidence) only when a margin gate and a
score-ratio floor pass; otherwise the verdict is \emph{attribution
abstain}. The v0.15 evidence mask records which decoded bits rest on actual
band evidence, so archived-uid cross-checks treat missing bands as
``unverifiable'' rather than ``contradicting''---eliminating a class of
false-abstain failures on no-attack round trips.

\subsection{Derived salts and metadata-loss recovery}
The default session salt is $\mathrm{HMAC}(K, \mathrm{sha256}(C(t)),
\mathrm{nonce})$ where $C(t)$ replaces in-group words by their group head:
embedding only swaps within groups, so $C$ is invariant between original
and watermarked text. Tracing enumerates nonce candidates (0--5 plus
head-truncation variants) and keeps the best score-to-threshold margin,
so a trace succeeds even if \emph{all} per-document metadata is lost. Older
random-salt documents without metadata remain untraceable (documented
limitation).

\subsection{Operational layer}
Adapters for OpenAI SDK, LangChain, LiteLLM proxy, AutoGen, and CrewAI
wrap outputs with fail-open middleware (any embedding error passes the
original text through); an HTTP service and CLI expose embed/trace; an
append-only audit log and Prometheus metrics cover operations. Key
rotation keeps old watermarks traceable via versioned keys.

\section{Experimental Setup}
\label{sec:setup}

\paragraph{Corpus (E-A1).}
Prior watermark evaluations often rely on synthetic or templated text,
which inflates detection rates: AI-generated text is dictionary-dense in
exactly the words watermark dictionaries cover. We instead generate
\textbf{1,020 real LLM deliverables} with Qwen2.5-32B via an API provider:
three genres (consulting report, legal memo, research survey section) in
two languages (Chinese 510, median 1,270 characters; English 510, median
3,738 characters), 24 seed topics, temperature 0.9, fixed generation seeds,
total cost under 2 CNY (usage logged per call). Generation prompts are in
the artifact.

\paragraph{Configuration.}
AAWM v0.15.0, zero-cost codec, bias 1.0, per-word vote cap 3, derived
salts. Two calibration configurations: uncalibrated defaults, and in-domain
calibration fitted on the first half of the corpus ($\ge$20 documents, per
the domain-calibration requirement established in our companion pitfalls
study). Tracing uses the explicit salt+bands metadata path; metadata-loss
recovery is a design feature evaluated separately in the artifact.

\paragraph{Metrics.}
Existence TPR (watermarked verdict on marked text), attribution exactness
(decoded user equals the embedding user), abstain rate, and under attack:
survival of both metrics. Embedding quality is measured as perplexity
change under a local Qwen2.5-3B (proxy metric; human evaluation is future
work).

\section{Detection and Attribution on Real LLM Deliverables}
\label{sec:robustness}

\subsection{Unedited leaks: the operating point}
Table~\ref{tab:ea0} gives the no-attack round trip on 100 held-out Chinese
deliverables with a single-user registry (uid 7): existence and attribution
are both 100\%, calibrated and uncalibrated. Real LLM deliverables sit
comfortably inside the watermark's capability window---their function-word
and connector vocabulary (``not only\ldots but also'', ``therefore'',
``should'') is precisely what zero-cost dictionaries cover.

\begin{table}[t]
\centering\small
\caption{E-A0: unedited-leak round trip on 100 real LLM Chinese
deliverables (zero-cost codec, explicit metadata path).}
\label{tab:ea0}
\begin{tabular}{@{}lccc@{}}
\toprule
Configuration & Existence TPR & Attribution exact & $n$ \\
\midrule
Uncalibrated defaults & 1.00 & 1.00 & 100 \\
In-domain calibrated & 1.00 & 1.00 & 100 \\
\bottomrule
\end{tabular}
\end{table}

\subsection{Real LLM paraphrase: the capability boundary}
We attack the marked texts with the same 32B LLM at three rewrite strengths
(``replace $\sim$15\%/30\%/50\% of content words, keep meaning and
length''; $n{=}50$ per strength, calibrated and uncalibrated).
Table~\ref{tab:ea2} shows the outcome: at 15\% rewrite strength, existence
TPR is 2\% and attribution is zero; at 30\% and 50\%, both are zero
everywhere. Real LLM paraphrase is far more destructive than the synthetic
within-group synonym attacks used in prior robustness reports, because the
rewriter varies exactly the connector and function words that carry the
mark, not just content words.

\begin{table}[t]
\centering\small
\caption{E-A2: survival under real LLM paraphrase (Qwen2.5-32B rewriter,
$n{=}50$ per cell; zero-cost Chinese codec). This is a capability boundary,
not a robustness table: semantic rewriting defeats lexical watermarks.}
\label{tab:ea2}
\begin{tabular}{@{}lccc@{}}
\toprule
Rewrite strength & Existence TPR & Attribution & Abstain \\
\midrule
15\% & 0.02 & 0.00 & 0.00 \\
30\% & 0.00 & 0.00 & 0.00 \\
50\% & 0.00 & 0.00 & 0.00 \\
\bottomrule
\end{tabular}
\end{table}

We therefore state the capability window plainly: AAWM traces
\emph{unedited or structurally edited} leaks (insertions, deletions,
reorderings, excerpting survive by design through content-addressed
voting). Against an attacker who semantically rewrites the leaked document
with an LLM, no lexical watermark survives; detection of the mark's
remains is a research problem, and the abstain protocol's honest output is
``watermark destroyed''. Deployments that expect rewrite attacks should
combine watermarking with traditional forensics rather than rely on any
lexical scheme.

\section{Embedding Quality}
\label{sec:quality}
Embedding rewrites every occurrence of a covered word toward its target
color---on the zero-cost dictionary, 34.5\% of covered tokens change
(substitution-rate measurement on the corpus). Whether this is perceptible
is an empirical question; we measure it with a local Qwen2.5-3B as a
perplexity proxy on 60 Chinese deliverables (sliding window, stride 512).
The token-level NLL ratio (marked/original) has mean 1.084, median 1.078,
and p95 1.138: a $\sim$9\% average perplexity increase, with no catastrophic
outlier in the tail. Before the v0.15.2 cross-boundary fragment guard, 4.0\% of
substitution sites (202/5070 on the corpus) were spurious
cross-word tokens (replacing \emph{意义} inside \emph{注意义务}
yields the garble \emph{注价值务}), which the guard now skips
conservatively; the PPL figures above are post-guard. We report this as
a \emph{proxy}: a 3B model's
sensitivity is weaker than a 7B model's or a human's, and a human blind
study is the mandatory next step before claiming ``imperceptible''.

\section{Discussion, Boundaries, and Related Work}
\label{sec:discussion}

\paragraph{Comparison.}
Against model-level watermarks (KGW and descendants; SynthID-Text), AAWM
trades sampling-level control for black-box applicability and per-user
attribution; against PostMark-class post-hoc watermarks, AAWM provides
multi-user uid decoding (16 bits with registry matching and abstention)
where presence-ratio schemes carry essentially one bit of existence;
against cryptographic agent identity, AAWM marks the text itself rather
than the agent's credentials. The honest price of uid capacity is
exposure to the pitfalls we characterize in the companion paper:
correlated-vote inflation (mitigated by per-word caps), domain-specific
null models (mitigated by in-domain calibration), and template-class
adversarial texts (mitigated by repetition routing)---all three mitigations
are shipped in v0.15.0.

\paragraph{Boundaries.}
(1)~Semantic LLM rewriting destroys the mark (Table~\ref{tab:ea2}).
(2)~The English zero-cost dictionary is sparse on non-AI text; production
English deployments should use the full-lexicon codec.
(3)~Legal-grade attribution is out of scope: the mark is statistical, not
cryptographic evidence.
(4)~A symmetric master key means verifier compromise is total; key
management is an operational precondition.

\paragraph{Artifact.}
AAWM is open source (MIT core) at v0.15.0: Python SDK, CLI, HTTP service,
five framework adapters, fail-open middleware, key rotation, audit
logging, and the full experiment suite of this paper with fixed seeds.


\bibliographystyle{plain}
\bibliography{refs}

\appendix
\section{Generation Prompts and Cost (E-A1)}
Prompts are genre instructions plus a topic drawn from a 24-topic pool
(manufacturing digitization, cross-border compliance, urban renewal,
EV supply chain, healthcare data governance, eldercare systems, energy
optimization, vocational education, carbon accounting, IP protection,
livestream regulation, personal-data protection, and 12 others per
language), temperature 0.9, max 1,400 tokens. Per-call usage and estimated
cost are logged in \texttt{api\_usage.jsonl}; total estimated cost for all
1,020 documents is 1.11 CNY (provider billing is authoritative).

\section{Embedding-Pipeline Reproduction of a Presence-Ratio Baseline}
As an external sanity check on statistic families, we re-ran PostMark's
official embedding pipeline (nomic word selection locally; insertion loop
faithful to the official v2 protocol: 10-word sublists, literal-presence
acceptance $\ge$0.5, $\le$3 attempts) with a 32B API inserter on C4 gold
texts. Detection-side results and the cross-environment embedding-table
nondeterminism we encountered are reported in the companion paper.


\end{document}
